\documentclass{article}
\usepackage{iclr2026_conference,times}
\usepackage{hyperref}
\usepackage{url}
\usepackage{graphicx}
\usepackage{booktabs}
\usepackage{amsmath}
\usepackage{xcolor}
\usepackage{amsfonts}

% -----------------------------------------------------------------------
% Working title — change before submission
% -----------------------------------------------------------------------
\title{Fine-Tuning Genealogy Is Written in Weights}

% Authors hidden for double-blind review.
% \iclrfinalcopy  % Uncomment for camera-ready only.

\newcommand{\fix}{\marginpar{FIX}}
\newcommand{\new}{\marginpar{NEW}}

% Placeholder macro for unwritten sections — shows in red, visible in draft
\newcommand{\TODO}[1]{\textcolor{red}{\textbf{[TODO: #1]}}}

% Shorthand for figure path prefix
\newcommand{\figpath}{figures}

\begin{document}
\maketitle

% -----------------------------------------------------------------------
\begin{abstract}
% TODO: draft after §3--5 are written.
% Template (from writing_prep.md):
%   Problem landscape → named formal contributions → breadth claim
%   (Flan-T5 full FT + LoRA variants + Llama-3.1-8B).
%
% Key numbers to hit:
%   WeightTree: 60\% sig at 7 leaves, 100\% at 8 (Fisher p < 0.0001)
%   LayerTrace: enc_sak (4\% of params) non-inferior to full weight
%   LoRA dose-response: Fisher-z r from -0.485 (k_only) to -0.805 (full FT)
%   Behavioral bridge: within-run r up to -0.40; null control r = -0.006
\end{abstract}

% -----------------------------------------------------------------------
\section{Introduction}
\label{sec:intro}
% -----------------------------------------------------------------------
% Arc (from writing_prep.md §Suggested intro arc):
%  1. Models proliferate, histories undocumented — the forking problem
%  2. Benchmarks tell you what a model does, not where it came from
%  3. Why it matters: SBOM analog, unlearning, adversarial consistency,
%     capability prediction
%  4. Two ways to probe lineage: behavioral (outputs) and whitebox (weights).
%     Each tells part of the story. Together they triangulate.
%  5. Research question (in italics)
%
% Key framing: "genealogical view" — not macro (embeddings for routing)
% or micro (pairwise fingerprinting), but full topology recovery against
% known ground truth. Name the monotone function claim.

\TODO{Write introduction.}


% -----------------------------------------------------------------------
\section{Related Work}
\label{sec:related}
% -----------------------------------------------------------------------
% Organize as: macro / micro / genealogical (our lane).
%
% Primary three to engage:
%   LLM DNA \citep{wu2026llmdna} — behavioral-space counterpart
%   Unsupervised MoTHer \citep{horwitz2025mother} — closest in goal, no RF
%   Neural Phylogeny \citep{yu2025neural} — pairwise edges, no full topology
%
% One-liner for all: "Prior work on [pairwise detection / behavioral
% phylogeny] lacks a controlled ground-truth evaluation — none can report
% RF distance against a known training topology."
%
% Also cite: PhyloLM, Neural Lineage, Nikolic et al., Zhu et al.
% Weight-space geometry: Model Soups, LMC, Task Arithmetic.

\TODO{Write related work.}


% -----------------------------------------------------------------------
\section{Framework}
\label{sec:framework}
% -----------------------------------------------------------------------

In this section, we define our terms and notation and establish the statistical framework underlying our experiments.


Much like \citep{wu2026llmdna}, we use the framework of phylogenies and LLM ``DNA'' to estimate important properties of parent-child training. We assume, unlike \cite{} but like \citep{yu2025neural}, that all node models are derived from a single root model. Unlike prior works, we focus not only on parent-child pairs but on recovering partial topology across all nodes, not just child-parent pairs.  We estimate the phylogeny by first estimating a distance matrix and then using neighbor-joining \citep{saitou1987neighbor} to recover the training tree.


\subsection{Phylogenetic Tree Recovery}
\label{sec:framework:phylo}
% We frame fine-tuning lineage recovery as a phylogenetics problem.
% Pairwise weight distances → NJ tree → RF vs ground truth → permutation p.
% Why NJ > directed MST (MoTHer): statistically principled, known properties.
% Cite: Saitou & Nei 1987 (NJ), Robinson & Foulds 1981 (RF).

To estimate $\widehat{\mathcal{P}}$, we use the following recipe:

\begin{enumerate}
    \item Choose set of layers $\mathcal{L}$ and distance (or similarity) metric ($d$)
    \item For layer $L^{(\ell)}$, estimate the distance between the layer weights from model $i$ to model $j$,
    $$
    D_{i,j}^{(\ell)} = d(L^{(\ell)}_i, L^{(\ell)}_j)
    $$
    \item Assemble $\hat{\textbf{D}}$, where $\hat{\textbf{D}}$ is function of the $D_{i,j}^{(\ell)}$ distances such as the average distance of pairwise weights across all layers
    \item Estimate $\widehat{{\mathcal{P}}}$ using $f(\hat{\textbf{D}})$ where $f$ is the NJ method\footnote{We assume that the models in question are leaves and not internal nodes}
\end{enumerate}

Step (3) allows us to use different layers to achieve different or even multiple estimates of $\widehat{\mathcal{P}}$, which we can do to evaluate layer importance.  Since phylogeny-space grows super-exponentially in the number of leaves, exact topology recovery is unlikely.
Instead, researchers commonly use Robinson-Foulds (RF) distance \citep{robinson1981comparison} to measure how ``close'' an estimated phylogeny is to the true one, a measure of partial topology recovery. An RF distance of $k$ means the two trees disagree on $k$ bipartitions, splits present in one tree but absent in the other.
A maximum RF distance can be computed from the total number of leaves, and the ratio (normalized RF) enables comparison across trees of different sizes. A normalized RF near 0 indicates near-perfect recovery; near 1 indicates a tree no better than random. We compare the RF of our estimated phylogeny against a null distribution of 1,000 uniformly random phylogenies with the same leaf count, yielding a permutation $p$-value per run.

%% that's t he statistical basis for our method.


\subsection{Distance Metric: Cosine and the SNR Argument}
\label{sec:framework:metric}

One question we sought to answer is what distance (or similarity) metric should we use to compare layers of weights to one another.  We decided to use vector-based norms as opposed to matrix or array based norms for easier portability across architectures.  The distances we considered are  cosine, correlation, $\ell_1$ and $\ell_2$ (the latter of which was used in MoTHER \cite{}).  The first two distances are bounded by 1 and so the distances are comparable across different layers.  The second two distances are unbounded, which can allow different magnitudes of weights to dominate the distance.

To determine which metric to use, we use within-run Signal to Noise Ratio (SNR), $\mathbb{E}[X]^2/\mathrm{Var}[X]$, as a heuristic with the intuition that a larger SNR will lead to easier separation among distances and thus better phylogenies.  The results of this experiment are shown in Table~\ref{tab:metric-snr}.

One result we see in Table~\ref{tab:metric-snr} is that cosine distance and correlation generally have the highest SNR across our different experiments.  Intuitively, this makes sense because the larger magnitude layers such as layer norms which are known to be noisy \cite{} have less chance to dominate.  Thus, we see that scaling the distance within layers is important.  For the rest of the analysis, we choose cosine similarity over correlation because we view them as largely as comparable with cosine distance having more prominence in layer weight analyses \cite{}.

\begin{table}[t]
\centering
\caption{Mean SNR ($= \mathbb{E}[X]^2/\mathrm{Var}[X]$ of per-layer RF values) with standard error in parentheses, by metric and tree size. Computed over \texttt{runs\_branching\_v3}. At 5--7 leaves all metrics are comparable; at 8 leaves normalized metrics (cosine, correlation) outperform $\ell_1$/$\ell_2$ by $\sim$18$\times$.}
\label{tab:metric-snr}
\begin{tabular}{lccccc}
\toprule
Metric & 5 leaves & 6 leaves & 7 leaves & 8 leaves & $n$ (5--7 / 8) \\
\midrule
Cosine      & 5.1 (4.0)   & \textbf{47.8} (24.7) & 60.4 (14.5) & 331 (270) & 9--15 / 4 \\
Correlation & \textbf{16.1} (8.6)  & 33.3 (13.5) & 71.7 (21.0) & \textbf{339} (269) & 11--15 / 4 \\
$\ell_1$    & 3.6 (2.1)   & 46.0 (27.4) & 76.2 (21.6) & 18.4 (8.0) & 9--13 / 3 \\
$\ell_2$    & 3.6 (2.1)   & 44.8 (27.4) & \textbf{81.3} (22.3) & 18.5 (7.9) & 9--13 / 3 \\
\bottomrule
\end{tabular}
\smallskip

\noindent\textit{Note:} Cell format is mean (SE). 4-leaf runs excluded (SNR $\approx 0$ across all metrics, $n$=2--3). Bold indicates highest value in each row and column. Large SEs at 8 leaves reflect high run-to-run variance with small $n$; the normalized vs.\ unnormalized gap is nonetheless $\sim$18$\times$.
\end{table}





\subsection{Per Layer Estimated Phylogenies}
\label{sec:framework:layers}
% Each weight tensor → its own NJ tree.
% Distribution of per-layer RF distances is the data; consensus tree is summary.
% More statistically principled than any competitor — flag explicitly.

One contribution we have in our methodology is the use of per-layer distance matrices between models as independent estimates as opposed to using the total distance across all layers.  The inspiration for this analysis is how biological phylogenies are estimated across multiple genes, the assumption being that different genes encode different information on how changes are occurring \cite{}.

Thus, a model with $L$ layers can result in $L$ distance matrices.  These distance matrices can either be used to each form their own estimate of $\hat{P}$ resulting in a collection of estimates.  From that set, a conensus phylogeny (i.e. the most common tree) is generally taken as the estimate.  Another alternative is to use a function of the distance matrices, such as the mean, to have a less noisy estimate of $\hat{\textbf{D}}$, to form $\widehat{\mathcal{P}}$. In addition to giving us a less variable estimate of the phylogeny, this layer-wise distance allows us to identify the most and least important layers in training by examining which layers change the most or least across multiple training runs.



% -----------------------------------------------------------------------
\section{Experiments}
\label{sec:experiments}
% -----------------------------------------------------------------------
Our methods bridge the gap between weight-level connections to peformance level connections such as predicted behavior.  We first create experiments to determine whether changes in weight structure can be used to recover (partial) the training trees, then determine whether distances in weights are associated with prediction capability in behavior.  We finally show an updated latent blackbox ancestry tree on some frontier models.

\subsection{Controlled Branching Trees}
\label{sec:exp:trees}
% Base: Flan-T5-base (250M). Bernoulli branching, 4 tasks, 50 replications,
% 5--8 leaves, known topology by construction.
% 2000 steps, lr=3e-4, batch=8, grad_accum=4, bf16.
% Controlled design enables RF evaluation — all competitors lack this.

To create a controlled experiment, we generate stochastic phylogenies or training trees, beginning from a root model.  The tree generating process is described in detail in Appenidx \ref{app:details}.  Our generating process produces fairly balanced phylogenies with up to 8 leaf nodes.  We generate 50 such phylogenies, each constituting the basic structure for a single training tree. We then assign task/dataset pairs to each edge. At each split in the tree, a child model is fine-tuned from its parent model on a random task/data set pair.  First a training task is randomly sampled and then a dataset within that task is randomly sampled.  The resulting pair is assigned to that edge.

%% This probably gets moved to intro or somewhere -- it's a real scenario
The motivation for this controlled experiment is from an industry-type training where a base model is given to different departments and each department is able to perform its own fine tuning on a model.  A question may arise if a new department begins and is interested in using an already fine-tuned model on a new task.  Which one should they pick?

The generated phylogenies are then used as fine-tuning instructions, the details of which can be found in Appendix [add-fine-tuning-instructions].  At each node in the phylogeny, a new model is trained and we can access each \textit{leaf} model's weights.  

Once we have the trained leaf models, we can use the steps in Section \ref{sec:framework:phylo} to estimate a phylogeny from the weight changes.  Since we know the actual training trees, we then can use RF distance to compare how far the estimated phylogeny is from the known training tree, unlike previous works.  We repeat this process for each generated phylogeny.

In our experiments, we vary the root model, task types, and which layers we train. We examine both an encoder-decoder model and a decoder-only model: Flan-T5-base (250M parameters) and Llama-3.1-8B, respectively. The tasks we fine-tune on include Q\&A, classification, translation, and text summarization. For Flan-T5-base, we run both full fine-tuning and four LoRA \citep{hu2022lora} adapter configurations varying the scope of trained parameters: key projection matrices only, query/key/value projections, all four attention projections (including output), and all projection matrices. Llama-3.1-8B serves as both a decoder-only and entry into more frontier-like models.  Llama-3.1-8B is trained with the key-only and query/key/value LoRA configurations; full fine-tuning is omitted at this scale due to memory constraints. Detailed training hyperparameters for all conditions are in Appendix~\ref{app:details}.

\subsection{Geometric Mechanism: Task Subspaces}
\label{sec:exp:geometry}
% Displacement subspace angles between fine-tuning directions.
% Cross-task grand mean ~84 degrees (near-orthogonal).
% Within-task 80.7--86.0 degrees — dose-response confirmed.
% RTE forgetting probe: classification vs RTE = 64.8 degrees.
% Cite: Ilharco et al. 2023 (Task Arithmetic) for task vectors.

To understand geometrically \textit{why} fine-tuning ancestry is recoverable from
weights, we measure the directions in which different training runs displace a model
from its starting point. For each fine-tuned checkpoint, we compute the weight
displacement $\Delta W = W_{\text{fine-tuned}} - W_{\text{base}}$, which captures
the net effect of all gradient updates relative to the pre-trained initialization.
We then measure the principal angles between the subspaces spanned by these
displacement vectors for pairs of fine-tuned models \citep{ilharco2023editing}.
Near-$90^\circ$ angles indicate that two training runs push the model into nearly
orthogonal regions of weight space; smaller angles indicate overlap.

We investigate three questions with this analysis. First, do different task types
induce near-orthogonal displacement directions. Second, does within-task training (same task type, different datasets) produce less orthogonal displacements than cross-task training? Third, does
LoRA's low-rank constraint alter the \textit{direction} of weight updates, or
does it primarily reduce their magnitude? We apply the subspace angle analysis
both globally (pooling all weight matrices) and per-layer, which can identify if certain layers are more variable in task/data set fine-tuning process.


\subsection{Behavioral Bridge}
\label{sec:exp:behavioral}
% Weight proximity predicts behavioral similarity (translation probe, semantic DV).
% Within-run r up to -0.40; no_translation_v2 null r = -0.006.
% Full regression details in Appendix~\ref{app:regression}.

While recovering ancestral information about a model is important for model provenance and auditing, it is also useful to have a practical interpretation of weight changes with regards to model behavior.  Here, we test whether difference in model weight distance is statistically significant in predicting model behavior on a new task.

We test for significance in two ways. First, we used a mixed effects regression model (see App. \ref{app:regression}), to determine whether the change in weight distance is a significant predictor for the outcome variable of change in benchmark score.  We also test for significance using the same model but for the outcome variable of semantic similarity from token embeddings.  

For Flan,  we train the model on trees where the hold-out task is translation.  For Llama, we train the model on the four standard tasks and the holdout is HellaSwag, a multiple choice data set where the model must guess the best sentence continuation.  The experiments differ in holdout task because we found that Flan performs terribly on HellaSwag, regardless of the fine-tuning task/dataset.



\subsection{Blackbox: Frontier Model Families}
\label{sec:exp:blackbox}
% GPT, Claude, Gemini, Llama probed via output embeddings across 4 tasks.
% Family structure recoverable when aggregating tasks.
% Single-task confounded by formatting (verbosity, markdown).
% Translation cleanest — language crossing decouples embedding from source content.

\TODO{Write blackbox experiment description.}


% -----------------------------------------------------------------------
\section{Results}
\label{sec:results}
% -----------------------------------------------------------------------

\subsection{WeightTree: Lineage Is Recoverable from Weights}
\label{sec:results:weighttree}

% Table 1: RF by leaf count
\begin{table}[t]
\caption{Tree recovery accuracy by tree size (branching\_v3, cosine distance).
  Observed RF is mean $\pm$ SE across runs; null mean from 1{,}000 random binary trees.
  \% sig: fraction of runs with $p < 0.05$ vs.\ random null.
  Fisher $p$: combined significance across runs.}
\label{tab:rf_by_leaves}
\begin{center}
\begin{tabular}{rrrrrr}
\toprule
Leaves & $n$ & Null mean (SE) & Observed RF (SE) & \% sig & Fisher $p$ \\
\midrule
5 & 11 & 3.21 (0.02) & 0.55 (0.21) &  0\% & 0.003 \\
6 & 12 & 5.24 (0.02) & 1.92 (0.36) & 50\% & 0.0003 \\
7 & 15 & 7.33 (0.01) & 2.60 (0.29) & 60\% & $<$0.0001 \\
8 &  4 & 9.42 (0.03) & 3.75 (0.48) & 100\% & $<$0.0001 \\
\bottomrule
\end{tabular}
\end{center}
\end{table}

% Figure 2: behavioral vs whitebox vs null
\begin{figure}[t]
  \begin{center}
    \includegraphics[width=\linewidth]{\figpath/behavioral_vs_whitebox.png}
  \end{center}
  \caption{RF distance by tree size for whitebox weight distances (cosine, blue triangles),
    behavioral distances (RTE label agreement, red squares), and random null (gray band,
    P10--P90). Whitebox sits well below the null at all tree sizes $\geq 5$ leaves;
    behavioral distances are indistinguishable from random. Behavioral failure is the
    point: it confirms that weight access is necessary for per-run lineage recovery.}
  \label{fig:behavioral_vs_whitebox}
\end{figure}

\TODO{Write WeightTree results prose.
  Headline: 60\% sig at 7 leaves, 100\% at 8.
  Per-run significance is the operational claim.
  Ref Figure~\ref{fig:behavioral_vs_whitebox} and Table~\ref{tab:rf_by_leaves}.}


\subsection{LayerTrace: Signal Concentrates in 4\% of the Model}
\label{sec:results:layertrace}

% Figure 3: per-layer RF profile
\begin{figure}[t]
  \begin{center}
    \includegraphics[width=\linewidth]{\figpath/layer_rf_profiles.png}
  \end{center}
  \caption{Per-layer RF distance (mean $\pm$ 2 SE, $n=47$ runs) for all five distance metrics.
    Encoder \textsc{sa.k} layers (highlighted) cluster at the bottom across all metrics ---
    consistently lowest RF, strongest lineage signal. The cyclical pattern
    (attention $\to$ FFN $\to$ layer norm repeating across blocks) shows that lineage
    concentrates in attention keys and dissipates through the FFN.
    % TODO: add Llama equivalent panel once Wright jobs complete
  }
  \label{fig:layer_rf_profiles}
\end{figure}

\TODO{Write LayerTrace results prose.
  enc\_sak (4\% of params) non-inferior to full weight at all tree sizes >= 5 leaves;
  outperforms at 6 leaves (RF 1.38 vs 1.92).
  Ref Figure~\ref{fig:layer_rf_profiles} and Table~\ref{tab:variants}.}


\subsection{Subspace Angles and Geometric Mechanism}
\label{sec:results:geometry}
% Cross-task ~84 degrees, FFN gates 88-89, SA.k/q 87-89.
% Within-task 80.7-86.0 -- lower than cross-task, dose-response confirmed.
% Layer norms pull within-task means down (52-80) while SA.k / FFN stay near-orthogonal.
% RTE probe: classification vs RTE = 64.8 degrees -- geometric forgetting mechanism.
% Subspace angle figure -> Appendix~\ref{app:robustness}.

\TODO{Write subspace angles results.}


\subsection{LoRA Dose-Response}
\label{sec:results:lora}

% Table 2: layer variants + LoRA variants combined
\begin{table}[t]
\caption{Lineage signal by training regime (branch structure analysis,
  cosine distance, \textsc{trained\_only=1} for LoRA groups).
  Fisher-$z$ $r$: Fisher $z$-averaged within-run correlation between
  cosine distance and same-branch indicator. \% sig: runs with $p < 0.05$.}
\label{tab:variants}
\begin{center}
\begin{tabular}{lrrrrr}
\toprule
Variant & Trained layers & Fisher-$z$ $r$ & SD & $N$ runs & \% sig (7-leaf) \\
\midrule
\multicolumn{6}{l}{\textit{LoRA adapter scope (Flan-T5-base)}} \\
k\_only        &  36 & $-0.485$ & 0.330 & 37 & --- \\
qkv            &  36 & $-0.594$ & 0.250 & 36 & 58\% \\
full\_attn     & 144 & $-0.671$ & 0.245 & 32 & 73\% \\
full\_ft\_approx & 168 & $-0.687$ & 0.242 & 31 & 70\% \\
\midrule
\multicolumn{6}{l}{\textit{Full fine-tuning}} \\
branching\_v3  & all & $-0.805$ & 0.230 & 35 & 60\% \\
\midrule
\multicolumn{6}{l}{\textit{LayerTrace variants (full FT, 6-leaf runs)}} \\
enc\_sak (4\%) & 12 tensors & \multicolumn{4}{l}{RF $1.38 \pm 0.18$ vs full $1.92 \pm 0.36$} \\
full weight    & 282 tensors & \multicolumn{4}{l}{RF $1.92 \pm 0.36$ (reference)} \\
\bottomrule
\end{tabular}
\end{center}
\end{table}

\TODO{Write LoRA dose-response prose.
  Key sentence: lineage survives LoRA at all.
  Large gap between full\_ft\_approx and full FT = constrained vs unconstrained updates.
  Ref Table~\ref{tab:variants}.}


\subsection{Behavioral Bridge}
\label{sec:results:behavioral}
% One paragraph. Full story in Appendix~\ref{app:regression}.
% Numbers: within-run r up to -0.40 (lora_k_only); all five groups significant.
% no_translation_v2 null: r = -0.006. Chain closed.

\TODO{Write behavioral bridge result (one paragraph).}


\subsection{Blackbox: Frontier Model Families}
\label{sec:results:blackbox}

% Figure 4: NJ tree grid (LLM DNA comparable)
\begin{figure}[t]
  \begin{center}
    \includegraphics[width=\linewidth]{\figpath/nj_tree_fingerprint_grid.png}
  \end{center}
  \caption{Neighbor-joining trees from behavioral fingerprint distances, one panel per task.
    GPT models (green) cluster consistently across all tasks.
    Claude clusters in summarization and QA but splits in translation (Haiku
    migrates due to its unique header-heavy formatting style in that task).
    Gemini--Llama proximity in QA and classification is a formatting artifact:
    both produce verbose, structured responses on those tasks, while GPT gives
    terse label-only answers. The artifact attenuates in summarization, where
    response length is task-constrained.}
  \label{fig:blackbox_nj}
\end{figure}

\TODO{Write blackbox results prose.
  Key points: Claude siblings tightest; Gemini-Llama proximity = formatting artifact;
  translation cleanest; multi-task aggregation necessary.
  Ref Figure~\ref{fig:blackbox_nj}.}


% -----------------------------------------------------------------------
\section{Discussion}
\label{sec:discussion}
% -----------------------------------------------------------------------

\subsection{The Provenance Chain}
\label{sec:discussion:chain}
% Closed chain: weights -> lineage (WeightTree), concentrated in attention keys
% (LayerTrace), survives LoRA (§results:lora), predicts behavior (§results:behavioral),
% visible in blackbox (§results:blackbox).
% Practical upshot: 4% of parameters sufficient; behavioral-only insufficient.

\TODO{Write provenance chain summary.}

\subsection{Why Attention Keys?}
\label{sec:discussion:keys}
% SA.k governs which token positions are attended to.
% Fine-tuning shifts routing toward task-relevant tokens; shift accumulates up tree.
% FFN stores task patterns in non-overlapping key-space regions (Geva et al. 2021).
% Encoder SA.k may be a universal "lineage layer" -- Llama results speak to this.

\TODO{Write mechanistic interpretation.}

\subsection{Whitebox vs.\ Blackbox: Complementary Roles}
\label{sec:discussion:wb_bb}
% Whitebox: exact genealogy, requires weight access.
% Blackbox: approximate family structure, API only.
% Behavioral projection story (§exp:geometry) explains why blackbox is lossy.

\TODO{Write whitebox vs blackbox discussion.}

\subsection{Limitations and Future Directions}
\label{sec:discussion:limitations}
% Controlled trees vs messy real-world genealogies.
% Single base model (Flan-T5) + Llama robustness check.
% LoRA rank threshold (post-ICLR rank experiment).
% Cross-architecture provenance = representational alignment problem.
% Prospective safety check: pre-training gradient angles predict forgetting.

\TODO{Write limitations and future work.}


% -----------------------------------------------------------------------
\subsubsection*{Acknowledgments}
\TODO{Add acknowledgments for camera-ready.}

% -----------------------------------------------------------------------
\bibliography{references}
\bibliographystyle{iclr2026_conference}


% -----------------------------------------------------------------------
\appendix
% -----------------------------------------------------------------------

\section{Experimental Details}
\label{app:details}

\subsection*{Tree Generation}
\label{app:tree-gen}

Training trees are generated via a Poisson branching process as follows:

\begin{enumerate}
    \item The root node (pretrained Flan-T5-base) draws $k \sim \mathrm{Poisson}(\lambda = 1.5)$ children, with a minimum of 1.
    \item Each subsequent node draws $k \sim \mathrm{Poisson}(\lambda = 1.5)$ children. Nodes that draw $k = 0$ naturally become leaves --- no separate pruning step is needed.
    \item Breadth First Search expansion continues until either a node budget ($n_\mathrm{nodes} = 14$ training nodes, not counting the root) is exhausted or a maximum depth ($n_\mathrm{layers} = 4$) is reached. The budget cap creates the truncation; any node alive when the budget runs out becomes a leaf.
    \item Each non-root node is assigned a task and dataset sampled without replacement from 4 task types (summarization, classification, QA, translation) and 36 datasets (9 summarization, 10 classification, 8 QA, 9 translation). The topology (parent--child structure) is fixed by a topology seed; dataset assignments vary across the 50 replications.
    \item Runs with fewer than 4 leaves are excluded from analysis, as trees with 3 or fewer leaves have trivial topology with no meaningful RF distance.
\end{enumerate}

This procedure yields trees with 4--8 observed leaves across the 50 replications, with variable depth and branching structure. The ground-truth topology is fully known by construction, enabling exact RF evaluation.

\subsection*{Training Configuration}
Base model: Flan-T5-base (250M parameters). Each node trained for 2,000 steps, learning rate $3 \times 10^{-4}$, batch size 8, gradient accumulation 4 (effective batch 32), warmup 200 steps, weight decay 0.01, bf16 precision. Up to 10,000 training examples and 1,000 evaluation examples per dataset, fixed dataset seed across all nodes to avoid data variation as a confound.

\section{Layer Variant Definitions}
\label{app:layers}
% enc_sak: encoder.block.{0-11}.layer.0.SelfAttention.k (12 tensors)
% encoder_h: enc_sak + encoder.block.{0-11}.layer.0.layer_norm (24 tensors)
% etc.

\TODO{Fill layer variant table.}

\section{Null Distribution Construction}
\label{app:null}
% 1000 random binary trees per run, same leaf count, uniform topology.
% p-value = fraction of null trees with RF <= observed RF.

\TODO{Fill null distribution description.}

\section{Blackbox Experiment Details}
\label{app:blackbox}
% Models, API versions, prompt templates, sentence-transformer embeddings.

\TODO{Fill blackbox details.}

\section{Metric Selection: SNR Analysis}
\label{app:snr}
% Ridge plots of SNR = E[X]^2 / Var[X] by metric family.
% Figures: snr_ridges.png (>=6 leaves), snr_ridges_all.png (all sizes).
\begin{figure}[h]
  \begin{center}
    \includegraphics[width=0.8\linewidth]{\figpath/snr_ridges.png}
  \end{center}
  \caption{SNR distributions by metric family ($\geq 6$-leaf runs).
    Normalized distances (cosine, correlation) dominate. See \S\ref{sec:framework:metric}.}
  \label{fig:snr_ridges}
\end{figure}

\TODO{Write SNR analysis text.}

\section{Behavioral Regression: Full Results}
\label{app:regression}
% Mixed-effects tables, HC3 sensitivity analysis, estimator comparison.
% Main text cites: within-run r up to -0.40, run FE + HC3 p=0.008.
% Report all estimators here; note that cluster-robust p=0.15 was inflated
% by between-run heteroscedasticity absorbed by run fixed effects.

\TODO{Fill regression appendix.}

\section{Within-Task and Robustness Results}
\label{app:robustness}
% Subspace angles by group, within-task branch structure,
% LoRA full RF profiles.
% Grand means: summarization_only 80.7, translation_only 82.4,
%              qa_only 85.8, classification_only 86.0 degrees.

\TODO{Fill robustness appendix.}

\section{Llama-3.1-8B Experiment Details}
\label{app:llama}
% Architecture notes, checkpoint filtering (MIN_CKPT=2), partial checkpoint issue.
% 30/30 runs per group; ~19-21/30 usable after min-leaves filter.

\TODO{Fill Llama appendix.}

\end{document}
